\section{Quadratic parity rigidity}\label{sec:parity}
This section removes a nonpivotal ambiguity without enumerating signed
tables. It is stated for arbitrary self-dual rank; its application to
rank five checks only tensor identities and mod-two matrices.

\begin{lemma}[Dimension-field spectral bound]\label{lem:spectral}
Let $R$ be a self-dual fusion ring, and let $\C$ be a hypothetical complex
categorification. Select the self-dual lifts in its sphericalization,
write their dimensions as $d_0=1,d_1,\ldots,d_n$, and set
\[
 K=\Q(d_1,\ldots,d_n),\qquad \delta=\sum_i d_i^2.
\]
Suppose a formal codegree $f$ of $R$ has minimal polynomial
$m_f(t)\equiv t^{\deg m_f}\pmod2$.
Then
\[
 [K:\Q]\le\operatorname{ord}_{t=0}(p_R(t)\bmod2).
\]
\end{lemma}
\begin{proof}
The lifts are self-dual: their dimensions are opposite and nonzero,
whereas spherical duality preserves dimension. Their signed quotient
has commuting real symmetric multiplication matrices $T_i$ and satisfies
\eqref{eq:signed}. It is commutative because every element of its basis
is self-dual and duality is an anti-automorphism. It is semisimple since
its real regular representation is closed under transpose.

Every formal codegree $f$ of $R$ belongs to $K$, and $\delta/f\in\OO_K$.
Indeed, inflate the corresponding character to the sphericalization.
Its formal codegree is $2f$ and the sphericalization has dimension
$2\delta$. By \cite[Theorem 2.13]{Ostrik}, a central simple has dimension
$2\delta/(2f)=\delta/f$. On forgetting, this is an integral linear
combination of the selected dimensions and their negatives. It is an
algebraic integer, and is nonzero, so $f\in K$ as well.

Put $e=[K:\Q]$. The Galois orbit of the dimension character has exactly
$e$ characters, and contributes the monic integral factor
\begin{equation}\label{eq:orbit-factor}
 Q_d(t)=\prod_{\sigma:K\hookrightarrow\CC}(t-\sigma(\delta))
\end{equation}
to the characteristic polynomial of $\sum_iT_i^2$.
Its degree is $e$, with character multiplicity, even when $\delta$
generates a proper subfield of $K$. Integrality follows from Galois
invariance and algebraic integrality of the eigenvalues; the complementary
Galois-stable factor is integral for the same reason.
In a normal overfield, reduce at any prime above two. All conjugates of
$f$ reduce to zero by the hypothesis on $m_f$. Integrality of every
conjugate of $\delta/f$ implies that all conjugates of $\delta$ do too.
Thus $Q_d(t)\equiv t^e\pmod2$.
Finally, $T_i\equiv N_i\pmod2$ gives
$t^e\mid\det(tI-\sum_iT_i^2)\bmod2=p_R(t)\bmod2$, as required.
\end{proof}

\begin{theorem}[Quadratic parity rigidity]\label{thm:parity}
Let $R$ be self-dual of rank $n+1$, with FP dimensions
\[
 s_0=1,\qquad s_i=A_i+\sqrt q\quad(1\le i\le n),
\]
where $q>1$ is squarefree, $q\equiv2$ or $3\pmod4$, and
$A_i\in\ZZ_{\ge0}$. Put $M=(N_{jj}^{\,i})_{1\le j,i\le n}$.
Suppose
\begin{equation}\label{eq:parity-hypotheses}
 1+\sum_i A_i^2=nq,\qquad
 \operatorname{ord}_{t=0}(p_R(t)\bmod2)=2,\qquad
 \ker_{\FF_2}M=\FF_2(1,\ldots,1).
\end{equation}
Assume some $A_j>0$ has irreducible fusion matrix $N_j$.
Then every complex categorification of $R$ admits a spherical structure
and has a pseudo-unitary Galois conjugate.
\end{theorem}
\begin{proof}
Use the notation of Lemma~\ref{lem:spectral}. The FP formal codegree is
\[
 D=1+\sum_i s_i^2=2nq+2\Bigl(\sum_iA_i\Bigr)\sqrt q.
\]
It generates $\Q(\sqrt q)$, and its monic minimal polynomial is
\[
 t^2-4nqt+4n^2q^2-4q\Bigl(\sum_i A_i\Bigr)^2\equiv t^2\pmod2.
\]
The lemma gives $[K:\Q]\le2$, while its proof gives $D\in K$.
Consequently $K=\Q(\sqrt q)$.

Since $\OO_K=\ZZ[\sqrt q]$, write $d_i=a_i+b_i\sqrt q$ with
$a_i,b_i\in\ZZ$. At least one $b_i$ is nonzero. Comparing irrational
coefficients in the signed square equations gives
\[
 2a_jb_j=\sum_iT_{jj}^{\,i}b_i.
\]
Let $g$ be the positive gcd of the nonzero $b_i$, and put $c_i=b_i/g$.
Divide before reducing modulo two: $Mc=0$ in $\FF_2^n$.
The integer vector $c$ is primitive, so its mod-two image is nonzero.
The kernel hypothesis forces all $c_i$ to be odd, hence all $b_i\ne0$.
This division is necessary; reducing $b$ directly could lose an even
common factor.

The characters $(1,d_i)_i$ and $(1,d_i')_i$, where prime denotes quadratic
conjugation, are distinct common eigenvectors of real symmetric matrices.
They are orthogonal, so
\begin{equation}\label{eq:orthogonality}
 1+\sum_i a_i^2=q\sum_i b_i^2.
\end{equation}
The spectral-radius inequality $|T_i|\le N_i$ bounds every conjugate
dimension by $s_i$. Therefore
\[
 |a_i|+|b_i|\sqrt q=\max(|d_i|,|d_i'|)\le A_i+\sqrt q.
\]
Since $|b_i|\ge1$, this implies $|a_i|\le A_i$. Thus
\[
 nq\le q\sum_i b_i^2=1+\sum_i a_i^2
       \le1+\sum_i A_i^2=nq.
\]
Equality holds throughout: $|b_i|=1$ and $|a_i|=A_i$.

Choose the embedding for which $|d_j|=s_j$ at an irreducible $N_j$ with
$A_j>0$, and set $v=|d|$ entrywise. From
$N_jv\ge|T_jd|=s_jv$, a positive left Perron vector forces equality.
Irreducibility then gives $v=s$, with the scalar fixed by $v_0=s_0=1$.
Put $\eta_i=\operatorname{sgn}(d_i)$, $\eta_0=1$, and
$U_{ik}^{\,\ell}=\eta_i\eta_k\eta_\ell T_{ik}^{\,\ell}$.
The vector $s$ is a positive character for both $U$ and $N$. Hence
\[
 0=\sum_\ell(N_{ik}^{\,\ell}-U_{ik}^{\,\ell})s_\ell.
\]
Each summand is nonnegative by \eqref{eq:signed}; therefore $U=N$.

Interchanging a selected lift with its opposite according to $\eta_i$
realizes this integer sign change. The selected and opposite
multiplicities in a product are $(N+U)/2$ and $(N-U)/2$. The latter
vanish, so the selected lifts form a fusion subcategory of the
sphericalization, equivalent to $\C$ under the forgetful functor.
The spherical structure restricts to it; compare
\cite[Section 5.1]{Ostrik}. Although one embedding was used in the
argument, $U=N$ is an equality of integer tensors and constructs the
section for the original category. Applying the chosen embedding to
its structure data gives positive dimensions $s_i$, hence the stated
pseudo-unitary Galois conjugate.
\end{proof}

\begin{remark}[Recognition from the tensor]
In rank five the vector $(A_i)$ need not be guessed. The irrational
coefficient of $s_j^2=\sum_iN_{jj}^i s_i$ gives
$2A_j=\sum_{i=1}^4N_{jj}^i$. Recover these four integers from the square
rows, set $q=(1+\sum_iA_i^2)/4$, and verify every dimension-product
identity. Squarefreeness, the mod-two characteristic-polynomial
multiplicity, the kernel and irreducibility are then exact finite tests.
The supplement implements this recognition for every census tensor;
no radicand or list of exceptional indices is supplied to the criterion.
\end{remark}
